% =============================================================================
% ch01_intro.tex — Chapter 1: Introduction | Biomimetic Navigator
% Compiler: XeLaTeX | Language: Hebrew primary, English technical terms via \en{}
% =============================================================================

\selectlanguage{hebrew}

\section{מבוא}
\label{sec:intro}

% ── 1.1 Background ──────────────────────────────────────────────────────────
\subsection{רקע ומוטיבציה}
\label{subsec:background}

הניווט האוטונומי של כלי-טיס בלתי-מאוישים (\en{UAV — Unmanned Aerial Vehicles})
בסביבות חסרות-\en{GPS} — כגון מבנים פגועים, מנהרות, יערות ורחובות צרים —
נותר אחד האתגרים הפתוחים המרכזיים בתחום הרובוטיקה האוטונומית.
שיטות ניווט מסורתיות המסתמכות על \en{GPS} מאבדות את יכולתן בפנים מבנים
ובסביבות עם הפרעות אלקטרומגנטיות, ואילו גישות \en{SLAM}
(\en{Simultaneous Localization and Mapping}) מבוססות-ראייה בלבד
סובלות מניוון (\en{drift}) ורגישות לתנאי תאורה. \cite{Thrun2005ProbRobotics}

הטבע, לעומת זאת, סיפק פתרונות אלגנטיים לבעיות אלה על פני מאות מיליוני שנות אבולוציה.
העטלפים (\en{Chiroptera}), בפרט, פיתחו מנגנון אקולוקציה
(\en{biosonar echolocation}) המאפשר להם לנווט, לצוד ולהתחמק ממכשולים
בחשכה מוחלטת בדיוק מילי-מטרי.
הם פולטים פולסי \en{FM} (\en{Frequency-Modulated}) בטווח 20–120 \en{kHz},
מנתחים את תבניות ההשתקפות בקליפת השמע (\en{auditory cortex}),
ובונים ייצוג פנימי של הסביבה בזמן-אמת. \cite{GriffithBatEcholocation}

עבודה זו מציגה את \textbf{NavigatorCrew} — פלטפורמת בינה מלאכותית
מרובת-סוכנים המבוססת על \en{CrewAI} \cite{CrewAIDocs},
שמטרתה לתרגם עקרונות ביולוגיים אלה לאלגוריתמי ניווט ביומימטי
לרחפנים קוואדרוטורים (\en{quadrotor UAVs}) בזמן-אמת.

% ── 1.2 The Biological Model: Bat Echolocation ──────────────────────────────
\subsection{המודל הביולוגי: אקולוקציה של עטלפים}
\label{subsec:bat_bio}

עטלפים ממשפחת \textit{Rhinolophidae} (עטלפי-פרסה) פולטים קריאות
\en{CF/FM} (גל-רציף / תדר-משתנה) דו-שלביות.
שלב ה-\en{CF} (\en{Constant Frequency}) מאפשר זיהוי מהירות יחסית
דרך תזוזת \en{Doppler}, ואילו שלב ה-\en{FM} מספק מידע מדויק על מרחק.
הטביעת האקוסטית המשוחזרת (\en{echogram}) מנותחת
על-ידי תאי-עצב מיוחדים ב-\en{inferior colliculus}
הרגישים לשילובים ספציפיים של עיכוב-זמן (\en{time-delay})
ותכולה תדרית (\en{spectral content}). \cite{Simmons1979BatSonar}

\figref{fig:bat_system} מציג סכמה של מסלול העיבוד הביולוגי מול
מסלול העיבוד המלאכותי שאנו מממשים:

\begin{figure}[H]
    \centering
    % PLACEHOLDER: Replace with fig_bat_vs_artificial.png from VisualizationEngineer
    \fbox{\parbox{0.9\columnwidth}{\centering
        \vspace{2.2cm}
        \textbf{[PLACEHOLDER — \en{Visio-AI} Figure]}\\[0.4em]
        \small דיאגרמת זרימה: מסלול עיבוד ביולוגי (עטלף) לעומת מלאכותי (רחפן)\\
        \small יוחלף על-ידי \en{fig\_bat\_vs\_artificial.png}\\
        \vspace{2.2cm}
    }}
    \caption{השוואת מסלולי העיבוד: אקולוקציה ביולוגית של עטלף (שמאל)
             לעומת הצינור המלאכותי של \en{NavigatorCrew} (ימין).
             \en{(Visio-AI diagram generated by VisualizationEngineer agent)}}
    \label{fig:bat_system}
\end{figure}

% ── 1.3 Problem Statement ───────────────────────────────────────────────────
\subsection{הגדרת הבעיה}
\label{subsec:problem}

נגדיר באופן פורמלי את בעיית הניווט הביומימטי.
תהי $\mathcal{E} \subset \mathbb{R}^3$ סביבה תלת-ממדית לא-ידועה,
ויהי רחפן $\mathcal{D}$ ממוקם במצב $\mathbf{x}_t = (p_t, \theta_t) \in SE(3)$
בזמן $t$, כאשר $p_t \in \mathbb{R}^3$ הוא וקטור המיקום
ו-$\theta_t \in SO(3)$ הוא מטריצת הסיבוב.

\textbf{בעיית ה-SLAM הביומימטית} היא לאמוד בו-זמנית:
\begin{enumerate}[noitemsep]
    \item את נתיב הרחפן $\mathcal{X} = \{\mathbf{x}_1, \ldots, \mathbf{x}_T\}$
    \item את מפת הסביבה $\mathcal{M}$ — ייצוג \en{occupancy voxel map}
    בצפיפות $r = 5\,\text{cm}$
\end{enumerate}
תוך שימוש יחידאי בחיישנים ביומימטיים $\mathcal{Z}_t = \{z_t^{\text{sonar}},\,
z_t^{\text{LiDAR}},\, z_t^{\text{vision}}\}$.

% ── 1.4 "Fancy Formula": FM Pulse Matched-Filter Correlation ────────────────
\subsection{נוסחת הבסיס: מסנן-התאמה לפולס FM}
\label{subsec:formula}

ניתוח פולסי ה-\en{FM} מבוסס על קורלציה צולבת
(\en{matched-filter cross-correlation}) בין האות הפלוט לאות המוחזר.
עבור פולס \en{LFM} (\en{Linear Frequency Modulated}) $s(t)$:

\begin{equation}
    s(t) = A\, \rect\!\left(\frac{t}{T}\right)
           \exp\!\left(j2\pi\left(f_0 t + \frac{B}{2T}\,t^2\right)\right),
    \quad 0 \le t \le T
    \label{eq:lfm_pulse}
\end{equation}

כאשר $A$ הוא האמפליטודה, $T$ משך הפולס, $f_0$ תדר ההתחלה,
ו-$B$ רוחב-הסרט (\en{bandwidth}).
מסנן-ההתאמה $h(t) = s^*(-t)$ מייצר פלט:

\begin{equation}
    y(\tau)
    = \int_{-\infty}^{\infty} r(t)\, h(\tau - t)\, dt
    = \int_{-\infty}^{\infty} r(t)\, s^*(t - \tau)\, dt
    \label{eq:matched_filter}
\end{equation}

שיא $|y(\tau)|$ מתקיים ב-$\tau = \tau_0 = 2d/c$,
כאשר $d$ הוא המרחק למכשול ו-$c = 343\,\text{m/s}$ מהירות הקול.
הרזולוציה הרדיאלית (\en{range resolution}) מתקבלת:

\begin{equation}
    \Delta r = \frac{c}{2B}
    \label{eq:range_resolution}
\end{equation}

עבור $B = 80\,\text{kHz}$ (ערך אופייני לעטלפים \textit{Eptesicus fuscus}),
מתקבל $\Delta r \approx 2.1\,\text{mm}$—
דיוק המספק קרוב לדיוק הביולוגי המתועד. \cite{Simmons1979BatSonar}

% ── 1.5 The Extended Kalman Filter for Sensor Fusion ────────────────────────
\subsection{נוסחת היתוך: מסנן קלמן מורחב}
\label{subsec:ekf}

שלב היתוך החישנים מבוסס על \en{EKF}
(\en{Extended Kalman Filter}) \cite{Kalman1960}.
שלב ה\en{predict}:

\begin{equation}
    \hat{\mathbf{x}}_{t|t-1} = f(\hat{\mathbf{x}}_{t-1|t-1},\, \mathbf{u}_t),
    \qquad
    P_{t|t-1} = F_t P_{t-1|t-1} F_t^{\top} + Q_t
    \label{eq:ekf_predict}
\end{equation}

ושלב ה\en{update}:

\begin{equation}
    K_t = P_{t|t-1} H_t^{\top}
          \bigl(H_t P_{t|t-1} H_t^{\top} + R_t\bigr)^{-1},
    \qquad
    \hat{\mathbf{x}}_{t|t} = \hat{\mathbf{x}}_{t|t-1}
                              + K_t\bigl(z_t - h(\hat{\mathbf{x}}_{t|t-1})\bigr)
    \label{eq:ekf_update}
\end{equation}

כאשר $F_t = \partial f/\partial \mathbf{x}\big|_{\hat{\mathbf{x}}}$
ו-$H_t = \partial h/\partial \mathbf{x}\big|_{\hat{\mathbf{x}}}$
הם היעקוביאנים של פונקציות המעבר והתצפית בהתאמה.
הסוכן \en{SensorFusionEngineer} יממש גרסה מורחבת הכוללת
אי-ודאות תלוית-גיאומטריה (\en{geometry-dependent noise covariance}). \cite{MurArtal2015ORB}

% ── 1.6 3D Trajectory Placeholder Figure ────────────────────────────────────
\subsection{נתיב תלת-ממדי לדוגמה}
\label{subsec:traj_placeholder}

\figref{fig:traj_3d} מציג נתיב תלת-ממדי אופייני הנוצר במהלך סימולציית
ניווט בסביבת מנהרה סינתטית.

\begin{figure}[H]
    \centering
    % PLACEHOLDER: Replace with fig_trajectory_3d.png from VisualizationEngineer
    \fbox{\parbox{0.9\columnwidth}{\centering
        \vspace{2.2cm}
        \textbf{[PLACEHOLDER — 3D Trajectory Plot]}\\[0.4em]
        \small נתיב רחפן תלת-ממדי — סימולציית מנהרה\\
        \small יוחלף על-ידי \en{fig\_trajectory\_3d.png}\\
        \small \en{matplotlib mpl\_toolkits.mplot3d | 300 DPI}\\
        \vspace{2.2cm}
    }}
    \caption{נתיב ניווט תלת-ממדי של הרחפן בסביבת מנהרה סינתטית.
             קו כחול — נתיב \en{Ground Truth};
             קו אדום — הערכת \en{SLAM} ביומימטי;
             נקודות אפורות — ענן-נקודות \en{LiDAR}.
             \en{(fig\_trajectory\_3d.png — generated by VisualizationEngineer)}}
    \label{fig:traj_3d}
\end{figure}

% ── 1.7 Paper Structure ──────────────────────────────────────────────────────
\subsection{מבנה המאמר}
\label{subsec:structure}

שאר המאמר מאורגן כך:
\begin{itemize}[noitemsep]
    \item \textbf{פרק~\ref{sec:bio_basis}}    — הבסיס הביולוגי: אקולוקציה ועיבוד שמע של עטלפים
    \item \textbf{פרק~\ref{sec:sensors}}      — מודאליות החישנים: \en{LiDAR}, סונאר, \en{Vision-AI}
    \item \textbf{פרק~\ref{sec:slam}}         — אלגוריתמי \en{SLAM} ביומימטיים
    \item \textbf{פרק~\ref{sec:fusion}}       — ארכיטקטורת היתוך החישנים
    \item \textbf{פרק~\ref{sec:algorithm}}    — האלגוריתם הביומימטי המוצע
    \item \textbf{פרק~\ref{sec:results}}      — תוצאות סימולציה וניתוח
    \item \textbf{פרק~\ref{sec:conclusion}}   — סיכום ועתיד
\end{itemize}
